\begin{anexosenv}
\partanexos
\renewcommand{\thechapter}{\Roman{chapter}}

\chapter{Corpo do e-mail de consulta à XP Investimentos e à MetaQuotes}
\label{anexo:email}

Este anexo apresenta o corpo do e-mail planejado para esclarecer, junto à corretora XP Investimentos e à fornecedora do MetaTrader 5 (MetaQuotes), o enquadramento do uso e da distribuição acadêmica dos dados de mercado, conforme discutido na Seção~\ref{sec:conformidade}.

\vspace{1em}
\noindent\textbf{Assunto:} Consulta sobre uso e distribuição acadêmica de dados de mercado (Mini-Índice Ibovespa) -- Projeto qData / UFC

\vspace{1em}
\noindent Prezados(as),

Meu nome é Luis Maia, sou estudante do Bacharelado em Engenharia de Software da Universidade Federal do Ceará (UFC), Campus de Russas, e integrante do Laboratório de Tecnologias Inovadoras (LTI). Escrevo para pedir esclarecimentos sobre o enquadramento do uso dos dados de mercado fornecidos por vocês no projeto descrito a seguir.

\noindent\textbf{Contexto do projeto}

No LTI, desenvolvo, como Trabalho de Conclusão de Curso, a plataforma qData. Ela centraliza dados financeiros históricos e em tempo real (\textit{market data}) para uso exclusivo de pesquisadores e estudantes do laboratório, em pesquisas de finanças quantitativas. O ativo estudado é o contrato futuro de Mini-Índice Ibovespa.

O fluxo atual é este:
\begin{enumerate}
    \item os dados são obtidos por meio do MetaTrader 5, conectado à conta em corretora [XP Investimentos];
    \item a plataforma armazena os dados em banco de dados de séries temporais;
    \item a plataforma os disponibiliza, via API autenticada (HTTPS e gRPC), aos pesquisadores cadastrados do LTI.
\end{enumerate}

Características do uso:
\begin{itemize}
    \item é estritamente interno, acadêmico e não comercial;
    \item não há venda, sublicenciamento nem exposição pública dos dados;
    \item o acesso é restrito a usuários cadastrados e autenticados;
    \item a plataforma não executa ordens nem faz recomendação de investimentos, sendo apenas infraestrutura de dados para pesquisa.
\end{itemize}

\noindent\textbf{Dúvida}

Sabemos que a B3 regula a distribuição e a redistribuição de \textit{market data} por meio de suas Políticas de Consumo e Comercial, que em regra exigem Termo de Adesão ou contrato próprio. Como os dados passam a ser consultados por mais de um pesquisador por meio da plataforma, gostaríamos de entender se isso pode ser considerado redistribuição. Pedimos, por isso, o esclarecimento formal dos seguintes pontos:
\begin{enumerate}
    \item O uso descrito é permitido pelos termos de uso da conta e do fornecimento de dados? Se não for, quais condições ou contratos seriam necessários? (Na mensagem à MetaQuotes: o uso do terminal MetaTrader 5 para extração programática de cotações, com posterior armazenamento e compartilhamento interno, é permitido pela licença?)
    \item Esse compartilhamento interno com pesquisadores configura redistribuição de dados de mercado nos termos da B3?
    \item Existe isenção, licença ou tratamento diferenciado para fins educacionais e de pesquisa acadêmica?
    \item Há limitações quanto ao armazenamento histórico dos dados, à sua granularidade e ao número de usuários com acesso?
    \item Caso seja necessário formalizar o uso, qual é o procedimento e quem é o contato responsável?
\end{enumerate}

Podemos enviar documentação adicional sobre a arquitetura do sistema ou participar de uma reunião, se isso ajudar. Um parecer por escrito seria muito valioso, pois será citado no trabalho acadêmico e orientará a adequação da plataforma.

Agradeço desde já a atenção e fico à disposição.

\vspace{1em}
\noindent Atenciosamente,

\noindent Luis Maia\\
Estudante de Engenharia de Software -- UFC Campus Russas\\
Laboratório de Tecnologias Inovadoras (LTI)\\
{[e-mail]} | {[telefone]}

\end{anexosenv}
